
Claim-Conditioned Probability (CCP) uses NLI-based conditioning to detect hallucinations, reporting +0.05--0.10 ROC-AUC improvements. We could not reproduce the baseline: claim-type mass ratio ($\rho_j$) values were 20-$80\times$ lower than the inferred range (median 0.0354 factual, 0.0103 creative vs inferred range 0.75--0.85), with no statistical separation ($p = 1.0$, Cohen's $d = -0.0635$). We tested whether CCP degrades on creative text (fiction, metaphor) versus factual text due to implicit factual-ontology assumptions. Root cause analysis revealed that DeBERTa-v3-base NLI, trained on SNLI/MNLI semantic similarity tasks, assigns approximately 90\% probability mass to the ''neutral'' class for factual verification tasks, mechanistically driving $\rho_j \to 0$. This is a case study of known NLI miscalibration issues for factual verification---SNLI/MNLI training objectives optimize for semantic similarity, not factual consistency. Methodologically, this failure teaches a critical lesson: measurement validity is prerequisite for hypothesis testing---when a metric produces values 20-$80\times$ outside the inferred range, you cannot distinguish ''hypothesis is wrong'' from ''measurement is broken.'' We adapt Dodge et al. (2019) reproducibility checklist for hallucination detection: (R1) report raw metric distributions, not just ROC-AUC; (R2) validate NLI calibration on known examples; (R3) document claim decomposition with inter-method agreement ($\alpha > 0.7$); (R4) provide public code with baseline replication notebooks. With 50+ hallucination detection papers published in 2024 citing NLI-based methods, transparent failures prevent costly replication waste across labs.

\textbf{Keywords}: hallucination detection, NLI domain adaptation, reproducibility, negative results, measurement validity
